\chapter[Multidimensional Reflected SDEs]{Multidimensional Reflected Stochastic Differential Equations}

\section{Normal reflection in a half-space}

Let
\[
  \R_+^d:=\R^{d-1}\times[0,\infty),
  \qquad
  e_d:=(0,\ldots,0,1).
\]

\begin{definition}[Half-space Skorokhod problem]
Let $f=(f^1,\ldots,f^d)\in C([0,\infty);\R^d)$ with
$f(0)\in\R_+^d$. A pair
$(g,\ell)\in C([0,\infty);\R^d)\times C([0,\infty);\R)$ solves the
normal-reflection Skorokhod problem for $f$ if
\begin{enumerate}[label=(\roman*)]
  \item $g(t)\in\R_+^d$ for every $t\geq0$;
  \item $g(t)=f(t)+e_d\ell(t)$;
  \item $\ell(0)=0$, $\ell$ is nondecreasing, and for every $t\geq0$,
  \[
    \int_0^t\1_{\{g(s)\notin\partial\R_+^d\}}\,\dd\ell(s)=0.
  \]
\end{enumerate}
\end{definition}

\begin{proposition}[Explicit half-space map]
The half-space problem has the unique solution
\begin{align*}
  g^i(t)&=f^i(t), &&i=1,\ldots,d-1,\\
  \ell(t)&=-\min_{0\leq s\leq t}\bigl(f^d(s)\wedge0\bigr),\\
  g^d(t)&=f^d(t)+\ell(t).
\end{align*}
Consequently, on every finite interval $[0,T]$, the reflected path and the
regulator satisfy the same Lipschitz and modulus estimates as the
one-dimensional Skorokhod map in the last coordinate.
\end{proposition}

\begin{proof}
Only the last coordinate is constrained. Apply the one-dimensional theorem
to $f^d$ and leave the first $d-1$ coordinates unchanged.
\end{proof}

Let $B$ be an $m$-dimensional Brownian motion and let
$a:[0,T]\times\R_+^d\to\R^d$ and
$\sigma:[0,T]\times\R_+^d\to\R^{d\times m}$ be measurable. A reflected SDE
in the half-space has the form
\begin{equation}
  X_t=X_0+\int_0^t a(s,X_s)\,\dd s
      +\int_0^t\sigma(s,X_s)\,\dd B_s+e_dL_t. \label{eq:halfspace-rsde}
\end{equation}
Here $X_t\in\R_+^d$, $L$ is continuous and nondecreasing with $L_0=0$, and
$\dd L$ is supported on $\{t:X_t\in\partial\R_+^d\}$. Global Lipschitz and
linear-growth assumptions on $a$ and $\sigma$ give existence and pathwise
uniqueness by the same fixed-point argument as in one dimension.

\section{Geometry of normal reflection in a domain}

Let $D\subset\R^d$ be a nonempty open connected set. The reflection direction
will point into $\overline D$.

\begin{definition}[Exterior-sphere normals]
For $x\in\partial D$ and $r>0$, define
\[
  \mathcal N_{x,r}
  :=\left\{n\in\R^d:|n|=1,
       B(x-rn,r)\cap D=\varnothing\right\},
  \qquad
  \mathcal N_x:=\bigcup_{r>0}\mathcal N_{x,r}.
\]
A vector in $\mathcal N_x$ is an inward unit normal. The domain satisfies a
uniform exterior-sphere condition if there is $r_0>0$ such that
$\mathcal N_x=\mathcal N_{x,r_0}\neq\varnothing$ for every
$x\in\partial D$.
\end{definition}

\begin{proposition}[Equivalent exterior-sphere inequality]
For $x\in\partial D$, $r>0$, and a unit vector $n$, the following are
equivalent:
\begin{enumerate}[label=(\roman*)]
  \item $n\in\mathcal N_{x,r}$;
  \item for every $y\in\overline D$,
  \[
    \langle y-x,n\rangle+\frac{|y-x|^2}{2r}\geq0.
  \]
\end{enumerate}
\end{proposition}

\begin{proof}
The exterior-ball condition is
$|y-(x-rn)|^2\geq r^2$ for all $y\in\overline D$. Expanding the square gives
exactly the inequality in (ii). Conversely, the inequality gives the closed
ball exclusion. If an interior point lay on the sphere, moving it slightly
toward the centre would produce an interior point inside the excluded ball,
so the open ball is disjoint from $D$.
\end{proof}

The exterior-sphere condition alone does not justify every uniqueness claim
for a nonconvex domain. A standard additional hypothesis prevents nearby
normal directions from becoming tangential to one another.

\begin{definition}[Uniform non-tangency condition]
The domain satisfies a uniform non-tangency condition if there are
$\delta>0$ and $\beta\geq1$ such that, for every $x\in\partial D$, one can
choose a unit vector $q_x$ satisfying
\[
  \langle q_x,n\rangle\geq\frac1\beta
\]
for every $y\in\partial D\cap B(x,\delta)$ and every
$n\in\mathcal N_y$.
\end{definition}

Bounded $C^2$ domains satisfy both geometric conditions, as do many convex
and piecewise smooth domains after the appropriate formulation at corners.
For nonsmooth domains the normal may be set valued, so the direction used by
the regulator must be recorded explicitly.

\begin{definition}[Normal Skorokhod problem]
Let $f\in C([0,T];\R^d)$ with $f(0)\in\overline D$. A triple
$(g,\ell,n)$ solves the normal Skorokhod problem if
\begin{enumerate}[label=(\roman*)]
  \item $g\in C([0,T];\overline D)$;
  \item $\ell$ is continuous, nondecreasing, and $\ell(0)=0$;
  \item $n$ is measurable with respect to the Stieltjes measure $\dd\ell$,
        $|n_s|=1$, and $n_s\in\mathcal N_{g_s}$ for
        $\dd\ell$-almost every $s$;
  \item
  \[
    g_t=f_t+\int_0^t n_s\,\dd\ell_s,
    \qquad
    \int_0^T\1_{\{g_s\in D\}}\,\dd\ell_s=0.
  \]
\end{enumerate}
\end{definition}

\begin{remark}
For a bounded $C^2$ domain, the inward normal is single valued and continuous
on the boundary. The normal Skorokhod problem is then well posed, and its map
is continuous on compact time intervals. In a general nonconvex or
nonsmooth domain, existence, uniqueness, and continuity require geometric
hypotheses such as the exterior-sphere and non-tangency conditions above;
one cannot infer pathwise uniqueness from the exterior-sphere condition
alone.
\end{remark}

\subsection{Penalization and its sign}

Suppose $D$ is smooth and a $C^2$ function $p:\R^d\to[0,\infty)$ satisfies
$p=0$ on $\overline D$ and, near $\overline D$,
$p(x)=\operatorname{dist}(x,\overline D)^2$. Outside the domain,
$\nabla p$ points outward. Hence a penalized approximation to inward normal
reflection has the form
\begin{equation}
  g^\varepsilon_t
  =f_t-\frac1{2\varepsilon}
       \int_0^t\nabla p(g^\varepsilon_s)\,\dd s, \label{eq:normal-penalty}
\end{equation}
not the same formula with a plus sign. The factor $1/2$ is conventional and
matches $\nabla\operatorname{dist}^2=2\operatorname{dist}\,n_{\mathrm{out}}$
where the nearest-point projection is smooth.

\begin{theorem}[Reflected SDE in a smooth domain]
Let $D$ be a bounded $C^2$ domain and let $n(x)$ be its inward unit normal.
Consider
\begin{equation}
  X_t=X_0+\int_0^t a(s,X_s)\,\dd s
      +\int_0^t\sigma(s,X_s)\,\dd B_s
      +\int_0^t n(X_s)\,\dd L_s. \label{eq:smooth-domain-rsde}
\end{equation}
If $a$ and $\sigma$ are globally Lipschitz in $x$, uniformly in $t$, and have
linear growth, then \eqref{eq:smooth-domain-rsde} has a unique strong solution
for every square-integrable $X_0\in\overline D$. If instead the coefficients
are bounded and continuous and the covariance is uniformly nondegenerate,
there is at least a weak solution.
\end{theorem}

\begin{proof}[Proof outline]
The deterministic normal Skorokhod map on a bounded $C^2$ domain is
nonanticipative and continuous. Euler approximations combined with tightness
give weak existence. Under Lipschitz coefficients, the deterministic
stability estimate and the usual stochastic estimates give pathwise
uniqueness. Weak existence plus pathwise uniqueness then yields a strong
solution.
\end{proof}

\section{General reflection fields}

Let $\mathcal K(x)\subset\R^d$ be a nonempty closed set of admissible inward unit
directions for each $x\in\partial D$. Normalizing directions is essential:
without it, the same finite-variation push can be represented by rescaling a
direction and the scalar regulator in infinitely many ways.

\begin{definition}[General Skorokhod problem]
Let $f\in C([0,T];\R^d)$ with $f(0)\in\overline D$. A pair $(g,k)$ solves the
Skorokhod problem for $(D,\mathcal K,f)$ if
\begin{enumerate}[label=(\roman*)]
  \item $g\in C([0,T];\overline D)$ and $k$ is continuous with finite
        variation and $k_0=0$;
  \item $g_t=f_t+k_t$;
  \item the variation measure $\dd|k|$ is supported on
        $\{t:g_t\in\partial D\}$;
  \item there is a $\dd|k|$-measurable process $\gamma$ such that
  \[
    k_t=\int_0^t\gamma_s\,\dd|k|_s,
    \qquad
    \gamma_s\in\mathcal K(g_s),\quad |\gamma_s|=1,
    \quad\dd|k|\text{-a.e.}
  \]
\end{enumerate}
If $(g,k)$ is unique, write
$\overline\Gamma(f)=(\Gamma(f),\mathcal R(f)):=(g,k)$ and call
$\overline\Gamma$ the extended Skorokhod map.
\end{definition}

\begin{proposition}[Nonanticipativity]
If the Skorokhod problem is unique on every finite interval, then the extended
map is nonanticipative: if $f$ and $\widetilde f$ agree on $[0,t]$, then their
solutions agree on $[0,t]$.
\end{proposition}

\begin{proof}
Restrict both solutions to $[0,t]$. They solve the same deterministic problem
there, so uniqueness gives equality. Continuity of the map is not needed for
this argument.
\end{proof}

Let $a:[0,T]\times\overline D\to\R^d$ and
$\sigma:[0,T]\times\overline D\to\R^{d\times m}$.

\begin{definition}[Strong reflected SDE]
On a prescribed stochastic basis carrying an $m$-dimensional Brownian motion
$B$, a strong solution is a pair $(X,K)$ of continuous adapted processes such
that $X_t\in\overline D$, $K$ has finite variation, and
\begin{equation}
  X_t=X_0+\int_0^t a(s,X_s)\,\dd s
      +\int_0^t\sigma(s,X_s)\,\dd B_s+K_t, \label{eq:general-rsde}
\end{equation}
where $K$ has a polar decomposition
$K_t=\int_0^t\gamma_s\,\dd|K|_s$ with
$\gamma_s\in\mathcal K(X_s)$ and $|\gamma_s|=1$ for $\dd|K|$-almost every $s$, and
$\dd|K|$ is supported on $\{s:X_s\in\partial D\}$.
\end{definition}

\begin{definition}[Weak reflected SDE]
A weak solution consists of a filtered probability space, an
$m$-dimensional Brownian motion on that space, an initial variable with the
prescribed law, and adapted processes $(X,K)$ satisfying all conditions in
the preceding definition. The stochastic basis and the Brownian motion are
part of the unknown.
\end{definition}

Whenever the extended Skorokhod map is defined,
\eqref{eq:general-rsde} is equivalently
\begin{align*}
  \eta_t
    &=X_0+\int_0^t a(s,\Gamma(\eta)_s)\,\dd s
      +\int_0^t\sigma(s,\Gamma(\eta)_s)\,\dd B_s,\\
  (X,K)&=\overline\Gamma(\eta).
\end{align*}

\section{Weak existence from a continuous Skorokhod map}

\begin{theorem}[Abstract weak existence]
Assume that the extended Skorokhod map is uniquely defined and continuous for
uniform convergence on compact time intervals. Suppose $a$ and $\sigma$ are
bounded and continuous, and let the initial law be supported on
$\overline D$. Then the reflected SDE \eqref{eq:general-rsde} has a weak
solution on every finite horizon.
\end{theorem}

\begin{proof}
Fix $T>0$ and a partition with mesh tending to zero. Let
$\kappa_n(s)$ be the left endpoint of the cell containing $s$. Construct the
Euler scheme
\begin{align*}
  \eta^n_t
    &=X_0^n+\int_0^t a\bigl(s,X^n_{\kappa_n(s)}\bigr)\,\dd s
      +\int_0^t\sigma\bigl(s,X^n_{\kappa_n(s)}\bigr)\,\dd B_s,\\
  (X^n,K^n)&=\overline\Gamma(\eta^n).
\end{align*}
Nonanticipativity makes the recursion well defined.

The drift parts are uniformly Lipschitz in time. For the martingale parts,
the Burkholder--Davis--Gundy inequality and boundedness of $\sigma$ give
\[
  \E|M^n_t-M^n_s|^4\leq C|t-s|^2.
\]
Kolmogorov's tightness criterion therefore makes $(\eta^n)$ tight in
$C([0,T];\R^d)$. Continuity of $\overline\Gamma$ makes
$(\eta^n,X^n,K^n)$ jointly tight.

Take a convergent subsequence. The vanishing mesh and continuity of the
coefficients imply convergence of the drift characteristics. Rather than
claiming convergence of stochastic integrals driven by changing Brownian
motions, identify the limit martingale $M$ through its conditional moments:
\begin{align*}
  M_t&:=\eta_t-X_0-\int_0^t a(s,X_s)\,\dd s,\\
  \langle M\rangle_t
    &=\int_0^t\sigma(s,X_s)\sigma(s,X_s)^\top\,\dd s.
\end{align*}
The discrete martingale identities pass to the limit, so $M$ is a continuous
martingale with the displayed bracket. The martingale representation theorem,
after enlarging the space if necessary, produces an $m$-dimensional Brownian
motion $\widetilde B$ such that
$M_t=\int_0^t\sigma(s,X_s)\,\dd\widetilde B_s$. Finally, continuity of the
Skorokhod map gives $(X,K)=\overline\Gamma(\eta)$. Thus the limit is a weak
solution.
\end{proof}

\begin{remark}[Why the martingale identification is needed]
A Skorokhod representation may place the $n$th approximation with a different
Brownian motion on a new probability space. Uniform convergence of those
Brownian motions and their integrands does not, by itself, imply
$L^2$ convergence of the stochastic integrals. The martingale and quadratic-variation identities provide the required stable
passage to the limit.
\end{remark}

\begin{theorem}[Continuous dependence]
Let $a_n,\sigma_n$ be uniformly bounded and converge to $a_0,\sigma_0$
locally uniformly on $[0,T]\times\overline D$. Let the initial laws
$\mu_n$ converge weakly to $\mu_0$, and assume that the same continuous
extended Skorokhod map is used for every $n$.

If the limit reflected SDE has uniqueness in law, then
\[
  X^n\Longrightarrow X^0
  \quad\text{in }C([0,T];\overline D).
\]
If, moreover, all equations are driven on one stochastic basis by the same
Brownian motion, the initial variables converge in probability, and the
limit equation has pathwise uniqueness, then
\[
  \sup_{0\leq t\leq T}|X^n_t-X^0_t|\xrightarrow{\P}0.
\]
\end{theorem}

\begin{proof}
The tightness estimates from the weak-existence proof are uniform in $n$.
Every subsequential limit solves the equation with coefficients
$(a_0,\sigma_0)$, by the same martingale-characteristic argument and the
continuity of $\overline\Gamma$. Uniqueness in law identifies every limit and
therefore proves weak convergence of the full sequence.

For the common-noise statement, apply the same argument to every pair of
subsequences. Any joint limit consists of two solutions of the limit equation
with the same Brownian motion and the same initial variable. Pathwise
uniqueness forces the two limits to coincide. The Gy\"ongy--Krylov criterion
then yields convergence in probability.
\end{proof}
